\section{Conclusion and Open Problems}\label{sec:conclusion}

Elkies' sum model holds up on a board without edges once the passed-pawn test respects en passant and the
frozen kings are checked. In the model, $\nim{2}$ appears with seven pawns and single-component $\nim{3}$ with nine
on ordinary boards and ten on the seven-file cylinder; with queen promotions and frozen kings the $1/8$ position
has oracle value exactly $1/8$. Six problems are open.

\begin{enumerate}
\item \emph{A certified $1/8$.} With queen promotions and frozen kings the oracle value is exactly $1/8$. Do
  the decisive end positions keep their results with all four promotions (Section~\ref{sec:real-eighth})?
\item \emph{The seven-pawn $\nim{2}$ in real chess.} The 8-man end positions deciding A and B agree with the
  model; those with 9 to 11 men are unsolved
  (Section~\ref{sec:real-eighth}).
\item \emph{Higher nimbers.} The least number of pawns in a ring $\nim{2}$ on $\cyl{n}$ is $n+3$ for $3\le n\le7$,
  ten for $7\le n\le10$ \lab{C2} and $n$ for every $n\ge10$ as a whole-board game \lab{S}. For $\nim{3}$ it is $n+2$ on $\cyl{10}$, $n+1$ on $\cyl{11}$ and at
  most $n+1$ up to $\cyl{59}$ \lab{C1}: can it be $n$? Are the wall
  families $\nim{2}$ as sums for all $n$? Is there a $\nim{4}$ on eight ranks?
\item \emph{When is a tr\'ebuchet zero?} Theorem~\ref{thm:cert} settles the promotion game, not chess: find a
  local test for chess, or settle Conjecture~\ref{conj:horizon2}; real tr\'ebuchets seem to have $h\approx5$--$7$ \lab{Conj}.
\item \emph{Endgames without edges.} With at most one new piece, $11\%$ of KPKP positions on the cylinder
  stay undecided, and KPPK comes next. Pawn play across the twisted seam of the M\"obius strip
  (Section~\ref{sec:tb}), where a pawn reverses direction, remains open.
\item \emph{Mis\`ere Dawson's chess.} Dawson asked about mis\`ere play \cite{Dawson1935}: is the mis\`ere
  quotient \cite{PlambeckSiegel2008} of Dawson's chess infinite? By the ring rule (Theorem~\ref{thm:ring}), an
  answer for strips also answers the cylinder.
\end{enumerate}
